% !PRINTED 614
\weberpage{614}
%<<CONT

5. \gesperrt{To every integral functional of the field $\Omega$ there belongs one and only one sequence of quadratic irrational numbers of the discriminant $\Delta$, and the numbers of this sequence differ by arbitrary integral rational numbers and are consequently equivalent to one another. Associated forms lead to the same sequence.}

If $\varphi_1$, $\varphi_2$ denote two integral functionals in $\Omega$, and we put
\[
\varphi_1 \varphi_2 = \varphi_3,
\]
then we can determine for each of these three forms a quadratic irrational number $\omega_1$, $\omega_2$, $\omega_3$, and $\omega_3$ is determined by $\omega_1$, $\omega_2$ up to an additive integral rational number. One then calls $\omega_3$ \gesperrt{compounded} or composed from $\omega_1$ and $\omega_2$, and sets, in symbolic notation,
\weq{3}{\omega_1 \omega_2 = \omega_3.}

If $\psi_1$ is equivalent to $\varphi_1$, and $\psi_2$ to $\varphi_2$, then by \S.\,152 $\psi_1 \psi_2 = \psi_3$ is also equivalent to $\varphi_3$.

Now if $\eta_1$, $\eta_2$, $\eta_3$ are the irrational numbers associated by Theorem 5. with the forms $\psi_1$, $\psi_2$, $\psi_3$, then by 3., \S.\,167 the numbers $\eta_1$, $\omega_1$; $\eta_2$, $\omega_2$; $\eta_3$, $\omega_3$ form three pairs of equivalent numbers, and we can therefore state the following theorem:

6. \gesperrt{If in a composition of quadratic irrational numbers each element is replaced by an equivalent number, then the compounded number also passes into an equivalent one.}

In this way the classes of the quadratic irrational numbers of a discriminant $\Delta$ are combined into an \gesperrt{Abelian} group isomorphic with the group of ideal classes.

In order actually to find the number compounded from two given numbers, one has to proceed as follows.

We shall denote the bases, formed according to (1), of the forms $\varphi_1$, $\varphi_2$, $\varphi_3$ as follows:
\[
\begin{array}{lll}
\varphi_1) & \alpha_1 = a_{1,1} & \alpha_2 = a_{1,2} + a_{2,2}\,\Theta \\
\varphi_2) & \beta_1 = b_{1,1} & \beta_2 = b_{1,2} + b_{2,2}\,\Theta \\
\varphi_3) & \gamma_1 = c_{1,1} & \gamma_2 = c_{1,2} + c_{2,2}\,\Theta,
\end{array}
\]
so that
\[
\omega_1 = \frac{\alpha_2}{\alpha_1}, \quad \omega_2 = \frac{\beta_2}{\beta_1}, \quad \omega_3 = \frac{\gamma_2}{\gamma_1}
\]
%CONT>>
